\subsection{Fixed-rank torus embeddings refute the unrestricted index hypotheses}

The next two findings concern the full subgroup restriction problem,
GCT VI, Problem 1.1.3, with the specifications of Section 3.4.
They do not concern its irreducible-inner plethysm subproblem,
Problem 1.1.2. The distinction is material: the broad problem permits
direct sums of polynomial torus characters.

\begin{proposition}[P-F14: fixed-rank modular indices are unbounded]
\label{prop:torusmodular}
For each integer $N\ge2$, let $H=\mathbb G_m$, $G=GL_2$ and
\[
 \rho_N(t)=\operatorname{diag}(t^{N+1},t),\qquad
 \lambda=(1),\qquad \pi=2.
\]
This is a faithful explicit polynomial embedding allowed by
GCT VI, Section 3.4.1. Its stretching multiplicity is
\[
 f_N(n)=\dim\operatorname{Hom}_{\mathbb G_m}
       (\mathbb C_{2n},\operatorname{Sym}^n(\mathbb C^2)|_{\rho_N})
       =\mathbf1_{\{N\mid n\}},\qquad n\ge0,
\]
and its generating function has minimum possible modular index $N$:
\[
 F_N(T)=\frac1{1-T^N}.
 \tag{P-F14}
\]
The input has $\operatorname{rank}(x)=5$ in the exact definition of
GCT VI, Section 3.4. Consequently Hypothesis 3.3.4 (PH3) is false for
the full class of explicit polynomial subgroup restrictions as stated.
Here $p(f_N)=s(f_N)=0$, so this example does not refute PH2 or SH(a).
\end{proposition}

\begin{proof}
The representation defining $\rho_N$ is the direct sum of the
polynomial characters of weights $N+1$ and $1$. Section 3.4.1,
``Basic explicit homomorphisms'', explicitly permits constructions
from irreducible polynomial representations by $\oplus$ and $\otimes$;
it specifies that polynomial representations of a torus are sums
of polynomial one-dimensional characters. Both exponents here are
positive. The map is faithful because its second diagonal entry is
$t$, and the inverse on its image is that entry. Its image is closed:
inside the diagonal torus it is given by $z_1=z_2^{N+1}$.

A monomial $x^k y^{n-k}$, $0\le k\le n$, in the symmetric power
has $H$-weight $(N+1)k+(n-k)=n+Nk$. It has weight $2n$ if and only
if $Nk=n$. There is exactly one such monomial when $N\mid n$ and
none otherwise; when it exists, $k=n/N$ belongs to $[0,n]$.
This proves the multiplicity and generating-function formulas.
Its natural graded multiplicity algebra is
$\mathbb C[xy^{N-1}]$, a polynomial ring with generator of degree $N$:
the weight equation forces every admissible monomial to be a unique
power of $xy^{N-1}$.

Let $\zeta$ be a primitive $N$th root of unity. The rational function
$F_N$ has a pole at $\zeta$. In any expression
\[
 F_N(T)=\frac{B(T)}{\prod_j(1-T^{a_j})},
 \qquad B\in\mathbb C[T],\quad a_j>0,
\]
the denominator must vanish at $\zeta$; otherwise the right side is
regular there. Thus $\zeta^{a_j}=1$ for some $j$, so $N\mid a_j$
and $\max_j a_j\ge N$. This argument does not need positivity of
$B$ or a bound on the number of denominator factors. The displayed
form (P-F14) attains the lower bound, proving minimality. Its
constituent polynomials are the constants zero and one, so both
indices are zero by GCT VI, Definitions 1.2.2 and 1.2.4.

Section 3.4 defines
$\operatorname{rank}(x)=\operatorname{rank}H+\operatorname{rank}G+
\|\lambda\|+\|\pi\|$.
Section 3.4.2 specifies $\|\lambda\|$ as the height for a
$GL_m$ partition, and specifies the combinatorial length of a torus
character as $\dim H$. Therefore the four terms here are
$1,2,1,1$, giving five independently of $N$. A polynomial in this
fixed parameter cannot bound the unbounded modular indices $N$.
This directly contradicts Hypothesis 3.3.4 in the stated broad scope.
\end{proof}

\begin{proposition}[P-F15: fixed-rank positivity and saturation indices are unbounded]
\label{prop:torusindices}
For odd $N\ge3$, let $H=\mathbb G_m$, $G=GL_3$ and
\[
 \sigma_N(t)=\operatorname{diag}(t,t^3,t^{N+1}),
 \qquad \lambda=(1),\qquad\pi=2.
\]
This is again a faithful explicit polynomial embedding. Its stretching
multiplicity $g_N$ satisfies
\[
 g_N(n)=\#\{(v,w)\in\mathbb Z_{\ge0}^2:2v+Nw=n\},
 \qquad
 \sum_{n\ge0}g_N(n)T^n=\frac1{(1-T^2)(1-T^N)}.
\]
In the definitions of GCT VI,
\[
 p(g_N)=s(g_N)=N-2,
 \qquad \operatorname{rank}(x)=6.
 \tag{P-F15}
\]
Its natural graded multiplicity algebra is a polynomial ring on two
generators of degrees $2,N$. Thus the ungraded affine scheme is the
same smooth normal affine plane for every $N$; over finite fields it
has pure Tate ordinary and intersection cohomology, as in P-F10.
Thus Hypothesis 3.3.2(a) (the uniform polynomial saturation-index
bound) and Hypothesis 3.3.3 (PH2) are false for the full subgroup
restriction class with the stated parameter. This finding does not
assert a refutation of the density assertion in Hypothesis 3.3.2(b).
\end{proposition}

\begin{proof}
The three characters have positive weights, so the same exact source
definition as in P-F14 applies. The first diagonal entry is $t$,
giving faithfulness and the inverse on the image. Inside the diagonal
torus, that image is closed by the equations
$z_2=z_1^3$ and $z_3=z_1^{N+1}$.
A monomial $x^u y^v z^w$ of degree $n$ has $u+v+w=n$ and weight
\[
 u+3v+(N+1)w=n+2v+Nw.
\]
It has the required weight $2n$ exactly when $2v+Nw=n$.
Conversely every nonnegative solution of that equation satisfies
$v+w\le n$, because both coefficients are at least two.
It therefore gives the unique admissible value $u=n-v-w$.
This is a bijection with the monomial weight basis and proves the
counting formula and its generating series.
In fact, $0\le w\le\lfloor n/N\rfloor$ and
$w\equiv n\pmod2$, because $N$ is odd. Counting those values gives
the closed formula
\[
 g_N(n)=\left\lfloor
 \frac{\lfloor n/N\rfloor+2-(n\bmod2)}2
 \right\rfloor.
\]
It holds also at $n=0,1$. Integer division and parity computation
evaluate it in time polynomial in the binary lengths of $n,N$.

More precisely, let $S=\mathbb C[x,y,z]$, with ordinary polynomial
degree and torus weights $1,3,N+1$. Its diagonal multiplicity algebra
\[
 T_N=\bigoplus_{n\ge0}(S_n)_{2n}
\]
is closed under multiplication, since degrees and weights both add.
The monomial condition is $u=v+(N-1)w$, so its monomial basis is
\[
 x^{v+(N-1)w}y^v z^w=(xy)^v(x^{N-1}z)^w,
 \qquad v,w\ge0.
\]
The exponent triples for distinct pairs $(v,w)$ are distinct.
Hence these two generators are algebraically independent and generate
all of $T_N$, proving
$T_N=\mathbb C[xy,x^{N-1}z]\simeq\mathbb C[U,V]$ with
$\deg U=2$, $\deg V=N$. The same monomial proof works over every
field. For $k=\mathbb F_q$ and any prime $\ell\nmid q$, P-F10's
affine-space computation therefore gives
$H_c^4(\operatorname{Spec}(T_N)_{\overline k},\mathbb Q_\ell)
=\mathbb Q_\ell(-2)$ and zero in all other degrees, with intersection
complex $\mathbb Q_\ell[2]$. This is an actual graded multiplicity
algebra, rather than only a ring with an equal Hilbert function.

P-F10, already proved in full, applies with $a=2$, $b=N$ since
$\gcd(2,N)=1$, and gives
$p(g_N)=s(g_N)=2N-2-N=N-2$. To make the exact index test explicit,
the constituents can also be read from
\[
 \frac1{(1-T^2)(1-T^N)}
 =\frac{\displaystyle\sum_{i=0}^{N-1}\sum_{j=0}^{1}T^{2i+Nj}}
 {(1-T^{2N})^2}.
\]
The exponents $r=2i+Nj$ represent all residues modulo $2N$ exactly
once: reduction modulo two determines $j$, and reduction modulo $N$
then determines $i$. The constituent belonging to $r$ is
$n\mapsto (n-r)/(2N)+1$. This is valid even when $0\le n<r$,
because the only possible such index in the residue is $r-2N$,
where the formula is zero. The greatest representative is $3N-2$.
After a nonnegative integer shift $c$, the smallest constant term
is $(c-(N-2))/(2N)$. All slopes are positive, so coefficientwise
positivity is equivalent to $c\ge N-2$. For strict saturation
each nonzero constituent must be positive at every integer $n\ge1$;
the smallest value at $n=1$ is $(1+c-(N-2))/(2N)$, which is
strictly positive precisely when $c\ge N-2$. This rechecks both
indices directly with the source conventions.

The four terms in the source's rank parameter are now $1,3,1,1$,
giving six. Since $N-2$ is unbounded, no polynomial in this fixed
parameter bounds either index. These are precisely the universal
bounds asserted in Hypotheses 3.3.2(a) and 3.3.3. A sequence of
counterexamples to a universal bound does not by itself determine
the density appearing in Hypothesis 3.3.2(b).
\end{proof}

\paragraph{Exact source scope and consequences.}
The primary locators are GCT VI, Problem 1.1.3; Hypotheses 3.3.1--3.3.4,
printed pp.~50--51; the definition of $\operatorname{rank}(x)$ in
Section 3.4, printed pp.~52--53; Section 3.4.1, printed pp.~53--54;
and Section 3.4.2, the specifications of $[\rho]$, items (4)--(5),
and of $[\lambda],[\pi]$, printed pp.~55--56.
The source expressly applies Hypotheses 3.3.1--3.3.4 to this
subgroup-restriction input. The maps above use its permitted
direct sums, rather than an irreducible representation of $H$.
Every irreducible representation of the torus is one-dimensional,
so these two- and three-dimensional representations are not inputs
to the irreducible-inner plethysm Problem 1.1.2. No refutation of
Hypotheses 1.6.6--1.6.7 for that narrower family follows.
The conditional implication PH3$\Rightarrow$PH2 in P-F12 is also
unaffected; it does not assert PH3 for these broader inputs.

Both examples even have homogeneous rational-polytope models of
constant ambient dimension. For P-F14 use the singleton
$\{1/N\}\subset\mathbb R$, whose $n$th dilation contains one
integer exactly when $N\mid n$. For P-F15 use
\[
 \{(v,w)\in\mathbb R_{\ge0}^2:2v+Nw=1\}.
\]
Their descriptions have binary length $O(\log N)$, as do the
torus-character inputs. Thus PH1-style efficient descriptions are
compatible with these failures of the index bounds. The indices
in question grow exponentially in this binary length; adding only
the bitlength of $\rho$ to the rank parameter would not make a
polynomial index bound true. A valid replacement must restrict
the family or retain suitable magnitude information about the
embedding weights. These explicit counterexamples do not identify
what occurs for the narrower determinant/permanent obstruction
families, and no historical priority claim is made.
Both examples have easy exact counting formulas---divisibility in
P-F14 and the displayed floor formula in P-F15---so their large
indices do not themselves imply a computational lower bound.

In particular, the broad rank-only bounds cannot follow solely from
purity, smoothness, normality, a positive monomial basis, or an
efficient PH1 polytope: the examples above have all these properties.
Quantitative control of the grading or embedding weights is a
separate requirement even when the graded algebra already realizes
the multiplicities exactly.
